\section{Excluding fourth and higher powers}
\label{sec:powers}

A spherical class whose homology suspension vanishes is a Pontryagin
square.  To apply the cone-square argument to such a class, we will need
its square root to have nonzero suspension.  At level zero this requires
a separate argument: a fourth power and its square root both disappear
under one homology suspension.  We exclude this possibility using the
weight projections and an integral characteristic class.

In this section, $*$ denotes the additive Pontryagin product.  A pairing
of spectra $E_1\wedge E_2\to E_3$ induces a based pairing of their zero
spaces, and we write $\circ$ for the corresponding pairing on homology.
Thus $*$ and $\circ$ are different operations.  We write $a^{*N}$ and
$a^{\circ N}$ for their respective powers whenever both are present.

\begin{proposition}
\label{prop:no-fourth-powers}
There is no positive-dimensional stable sphere class whose nonzero
mod-$2$ Hurewicz image is a fourth or higher iterated Pontryagin
square.  Equivalently, if $f:S^d\to Q_0S^0$, $d>0$, and
\[
 h(f)=z^{*2^a},\qquad a\geq2,
\]
then $h(f)=0$.
\end{proposition}

Eccles and Zare exclude higher iterated squares when the root survives
homology suspension \cite[Theorem~3.10]{EcclesZare2009}, and state that
the root of a spherical square is odd-dimensional
\cite[Lemma~3.12]{EcclesZare2009}. Here we prove the required exclusion
using spectrum pairings and integral characteristic classes.

The space maps adjoint to individual weight projections need not
preserve Pontryagin squares. We first compute their correction term,
which is expressed through the block pairings of extended powers.

For $v\geq1$, set
\[
 X_v=(B\Sigma_v)_+,\qquad E_v=D_vS\simeq\Sigma^\infty X_v.
\]
Block inclusion gives pairings $E_v\wedge E_w\to E_{v+w}$, associative
and commutative coherently, and thus associative and commutative
pairings $\circ$ on their zero-space homology.  Write
\[
 \varepsilon_v:\Sigma^\infty QX_v\longrightarrow E_v
\]
for the stable counit; its induced map on positive-dimensional homology
is stable homology evaluation.

Fix an integer $N>0$.  Let
\[
 \tau:\Sigma^\infty_+QS^0\longrightarrow\mathbb P_{\leq N}(S\{u\})
\]
be the augmented commutative ring projection of
\eqref{eq:weight-projections}.  For $1\leq i\leq N$, let
\[
 j_i:QS^0\longrightarrow\Omega^\infty E_i=QX_i
\]
be the based adjoint of its weight-$i$ projection, and write
$J_i=(j_i)_*$ on positive-dimensional homology.  These are generally
\emph{not} additive infinite loop maps.  Their multiplication formula
is the following.

\begin{lemma}[Multiplication under the weight projections]
\label{lem:weight-product}
For points $x,y\in QS^0$, there is a homotopy
\begin{equation}
\label{eq:weight-product}
 j_i(x+y)\simeq j_i(x)+j_i(y)
       +\sum_{0<t<i}j_t(x)\circ j_{i-t}(y),
\end{equation}
where addition in the target is its additive infinite loop structure.
For positive primitives $A_1,A_2\in H_*Q_0S^0$, this implies
\begin{equation}
\label{eq:primitive-weight-product}
 J_i(A_1*A_2)
 =(J_iA_1)*(J_iA_2)
   +\sum_{0<t<i}(J_tA_1)\circ(J_{i-t}A_2).
\end{equation}
Each $J_i$ preserves primitivity and $A$-annihilation.  In particular,
\begin{equation}
\label{eq:weight-square}
 J_i(A^{*2})=(J_iA)^{*2}
 +\begin{cases}(J_{i/2}A)^{\circ2},&i\text{ even},\\0,&i\text{ odd}.
 \end{cases}
\end{equation}
\end{lemma}

\begin{proof}
Under the adjunction between suspension ring spectra and
$E_\infty$ spaces, $\tau$ gives an $E_\infty$ map from the additive
structure of $QS^0$ to the multiplicative space of the truncated ring.
The source is grouplike, so this map lands in units.  Its weight-zero
coordinate is constantly $1$, since $\tau$ is augmented; the
positive coordinates at the source basepoint are zero.  Multiplication
in the truncated graded ring has weight-$i$ coordinate equal to the
right side of \eqref{eq:weight-product}.  This proves that formula,
including its compatibility with all lower coordinates in the single
chosen truncation.  Group completion causes no further terms:
$1+u$ is a unit in the target, so the map is precisely the extension
of the free-monoid map through the localization
\eqref{eq:zero-localization}.

To evaluate \eqref{eq:weight-product} on $A_1\times A_2$, distribute
the two primitive diagonals across its additive summands.  Either the
two positive inputs enter the two linear terms, giving their
Pontryagin product, or both enter one bilinear summand, giving its
block pairing.  Every other allocation pairs a positive class with
the zero point and vanishes.  This gives
\eqref{eq:primitive-weight-product}.  Although $j_i$ need not preserve
addition, it is a based space map, so it preserves ordinary
primitivity and commutes with the dual Steenrod operations.
When $A_1=A_2$, the terms with indices $t$ and $i-t$ cancel by block
commutativity; only the equal-block term can remain.  This proves
\eqref{eq:weight-square}.
\end{proof}

To remove the extra block-square term in \eqref{eq:weight-square},
we will use nilpotence for primitives whose stable homology evaluation
vanishes. The block pairings arise from spectrum maps, so they have
coherent additive structure in each variable. The following calculation
uses that structure; it does not require the maps $j_i$ to be additive.

\begin{lemma}[Dyer--Lashof operations and spectrum pairings]
\label{lem:pairing}
Let $E_1\wedge E_2\to E_3$ be a pairing of spectra.  If
$u\in H_a\Omega^\infty E_1$ is homogeneous and
$y\in H_b\Omega^\infty E_2$, where $b>0$, is $A$-annihilated, then,
for $i\geq a$,
\begin{equation}
\label{eq:pairing-rule}
 (Q^i u)\circ y=
 \begin{cases}
 Q^i(u\circ y),&i\geq a+b,\\
 0,&a\leq i<a+b.
 \end{cases}
\end{equation}
If, in addition, $y$ is primitive and belongs to the basepoint component,
then $(u_1*u_2)\circ y=0$ for positive-dimensional $u_1,u_2$.
Translating a positive-dimensional first factor by a component point
does not change its pairing with $y$.
\end{lemma}

\begin{proof}
Put $X=\Omega^\infty E_1$ and $Y=\Omega^\infty E_2$.
The pairing is a coherent family of additive infinite loop maps in its
first variable.  Indeed its adjoint is a spectrum map
$E_2\to F(E_1,E_3)$; passing to zero spaces gives a map from $Y$ to the
space of spectrum maps $E_1\to E_3$, and hence to the space of additive
infinite loop maps.  Connective covers may be taken without changing
the zero spaces.  In particular, the naturality homotopies for the
additive $E_\infty$ structure vary coherently over $Y$.

For an unbased space $U$, write
\[
 \mathcal E_2U=EC_2\times_{C_2}(U\times U),
\]
where $C_2$ interchanges the factors.  The naturality just described
factors the pairing of an additive square operation with $Y$ through
\[
 \delta:\mathcal E_2X\times Y\longrightarrow\mathcal E_2(X\times Y),
 \qquad [e;x_1,x_2],y\longmapsto[e;(x_1,y),(x_2,y)].
\]
Let $W$ be the free periodic $\F_2[C_2]$-resolution with generators
$e_j$ and differential $d e_j=(1+\tau)e_{j-1}$.
The chain model
\[
 W\otimes_{C_2}C_*(U)^{\otimes2}
\]
computes $H_*\mathcal E_2U$.  Choosing cycles for a homology basis of
$U$ gives an equivariant quasi-isomorphism before taking homotopy
orbits.  It follows that a basis of the resulting homology consists of
unordered unequal pairs in resolution degree zero and classes
$e_j\otimes z\otimes z$, $j\geq0$, on equal basis elements.

We claim the universal formula
\begin{equation}
\label{eq:external-pairing}
 \delta_*\bigl((e_j\otimes u\otimes u)\times y\bigr)
 =\sum_{t\geq0}e_{j-b+2t}\otimes
       (u\times\Sq_*^t y)\otimes(u\times\Sq_*^t y),
\end{equation}
with terms of negative resolution degree omitted.  This is the
extended-square calculation underlying May's formula
\cite[\S2, formula~(2.11), pp.~246--247]{Madsen1975}; we include its verification.

Transfer from $\mathcal E_2U$ to $U^2$ sends an unequal pair to the sum
of its two ordered versions and kills all equal-pair classes.  Thus it
detects the unequal-pair part injectively.  Transfer is natural for
$\delta$, since its pulled-back double cover is the double cover of
$\mathcal E_2X\times Y$ obtained from $\mathcal E_2X$.
The input of \eqref{eq:external-pairing} has transfer zero.  The output
therefore has no unequal-pair terms.

For $\beta\in H^qU$, let $P(\beta)\in H^{2q}\mathcal E_2U$ be its
external square, and let $h\in H^1BC_2$ be the characteristic class of
the double cover.  In the dual resolution model, a cocycle $f$
representing $\beta$ gives $P(\beta)$ by $f\otimes f$ in resolution
column zero.  This description is independent of $f$: if
$f'=f+dc$, then $P(f')-P(f)$ is the total coboundary of the cochain
whose columns zero and one are, respectively,
\[
 f\otimes c+c\otimes f+c\otimes dc,
 \qquad c\otimes c.
\]
Multiplication by $h^l$ is the resolution shift, so
$h^lP(\beta)$ is represented by $f\otimes f$ in column $l$.
It evaluates on $e_j\otimes z\otimes z$ as
$\delta_{lj}\langle\beta,z\rangle^2$.

External squares respect cross products over the common base $BC_2$.
Moreover, pulling $P(\beta)$ back along the Borel diagonal
$BC_2\times U\to\mathcal E_2U$ gives
\[
 \sum_{t=0}^{q}h^{q-t}\Sq^t\beta;
\]
this is the total-square definition of the Steenrod operations.
Consequently, for $\alpha\in H^*X$ and $\beta\in H^qY$,
\[
 \delta^*\bigl(h^lP(\alpha\times\beta)\bigr)
 =h^lP(\alpha)\sum_{t=0}^{q}h^{q-t}\Sq^t\beta.
\]
Its evaluation on the input in \eqref{eq:external-pairing} can be
nonzero only when
\[
 |\alpha|=a,\qquad q+t=b,\qquad j=l+q-t=l+b-2t.
\]
The coefficient is then
$\langle\alpha,u\rangle\langle\beta,\Sq_*^t y\rangle$.
Tensor cohomology classes detect all equal-pair coefficients; in each
evaluation only finitely many homology basis elements are involved.
Together with the transfer calculation, these evaluations prove
\eqref{eq:external-pairing}.

Now take $j=i-a$ and apply the label pairing, followed by the additive
square operation.  The term with index $t$ becomes
$Q^{i+t}(u\circ\Sq_*^t y)$.  When $y$ is $A$-annihilated, only $t=0$
remains.  Its resolution index is $i-a-b$, which proves both cases of
\eqref{eq:pairing-rule}.

Finally, for a primitive $y$ in the basepoint component,
$\Delta y=y\otimes[0]+[0]\otimes y$.  Distributivity of a spectrum
pairing over addition gives
\[
 (u_1*u_2)\circ y
 =(u_1\circ y)*(u_2\circ[0])
  +(u_1\circ[0])*(u_2\circ y)=0,
\]
because pairing a positive-dimensional class with the zero point is
zero.  If one first factor is a component point $[c]$, its pairing
with $[0]$ is the additive unit $[0]$ in the target; the same formula
proves the translation assertion.
\end{proof}

\begin{lemma}[Triple block nilpotence]
\label{lem:block-nilpotence}
Suppose $B\in H_i(Q_0X_v)$, $i>0$, is primitive and $A$-annihilated.
If $\varepsilon_{v*}B=0$, then
\[
 B^{\circ3}=0.
\]
\end{lemma}

\begin{proof}
The free infinite loop space homology theorem and group completion
describe $H_*Q_0X_v$ as a polynomial Pontryagin algebra
\cite[Part I]{CLM1976}.  Its generators are the positive-dimensional
particle classes from a basis of $H_*B\Sigma_v$, and the admissible
positive-excess Dyer--Lashof words of positive length on particle
classes, including the degree-zero particle.  Each is translated to
the basepoint component.  There is one degree-zero particle component
to invert in group completion; these translations do not change
stable homology evaluation in positive dimensions.

The counit takes the length-zero positive-dimensional generators
independently to their corresponding basis classes in $H_*E_v$.
Stable homology evaluation kills products of positive-dimensional classes.
It also kills positive-length words: upper-index operations commute with
homology suspension, and repeated suspension eventually makes their
outer operation unstable.  Therefore, modulo Pontryagin decomposables,
the hypothesis $\varepsilon_{v*}B=0$ expresses $B$ as a sum of
positive-length words only.

Pair this expression with $B$.  The preceding lemma discards the
decomposables and translations in the first factor.  A remaining term
has the form $Q^q c$ with $q+|c|=i$.  Its pairing can be nonzero only if
\[
 q\geq |c|+i=q+2|c|,
\]
so $|c|=0$.  In this case
\[
 (Q^i c)\circ B=Q^i(c\circ B)=(c\circ B)^{*2}.
\]
Thus $B\circ B$ is Pontryagin decomposable, and another pairing with
$B$ is zero by Lemma~\ref{lem:pairing}.

We also record a property used below: the block pairing of two
positive primitives in their basepoint components is primitive.
Indeed, naturality of the ordinary diagonal leaves the two end terms;
each crossed term pairs a positive-dimensional class with the zero
point and vanishes.
\end{proof}

\begin{proof}[Proof of Proposition~\ref{prop:no-fourth-powers}]
It suffices to exclude fourth powers, since every higher iterated square
is a fourth power. Suppose, therefore, that
\[
 h(f)=z^{*4}\ne0.
\]
Applying Lemma~\ref{lem:primitive-roots} twice shows that the
homogeneous positive-dimensional root \(z\) is primitive and
\(A\)-annihilated.

\emph{Choosing the detector.}
Let \(r=2^{k_0}\ge2\) be the leading weight of this root, as in
Lemma~\ref{lem:zero-weights}, and let \(z_r\in H_*D_rS\) be its
leading component. Put
\[
 x=z^{*2},\qquad b=|x|,\qquad s=2r=2^\ell.
\]
Here \(b\) is positive and even, \(\ell\ge2\), and \(d=2b\).
In the polynomial associated weight algebra the leading component of
\(x\) is
\[
 p=z_r^2\in H_bD_sS,
\]
and is nonzero. This square is taken in the graded homology algebra
\(H_*\mathbb P(S)\), rather than by the unstable pairing \(\circ\).

The class $p$ is primitive for the linear weight coproduct and is a
product of positive-dimensional factors.  It therefore dies under
the weight suspension \eqref{eq:weight-suspension}.
One may see this before group completion: homology suspension from
the free monoid kills positive-dimensional Pontryagin products, and
Lemma~\ref{lem:weight-suspension} identifies its weight components.
Thus, under the dual description in
Lemma~\ref{lem:weight-modules}, the functional defined by $p$ on the
Dickson algebra $D_\ell$ kills $d_0D_\ell$.
It is nonzero on some monomial in $d_1,\ldots,d_{\ell-1}$.
The restriction description of that lemma lifts $d_j$ to
$w_{s-2^j}(\rho_s)$, where $\rho_s$ is the permutation representation
of $\Sigma_s$.  Hence there is a class
\begin{equation}
\label{eq:even-characteristic-detector}
 v\in H^b(B\Sigma_s),\qquad \langle v,p\rangle=1,
\end{equation}
which is a polynomial in the \emph{even-index}
Stiefel--Whitney classes of $\rho_s$.

\emph{Removing the block correction.}
Use one truncation through weight \(s=2r\) to define all the maps
\(j_i\). We claim
\begin{equation}
\label{eq:power-projection}
 \begin{split}
 J_s h(f)
  &=(J_sx)^{*2}+(J_rx)^{\circ2}\\
  &=(J_sx)^{*2}+(J_{r/2}z)^{\circ4}\\
  &=(J_sx)^{*2}.
 \end{split}
\end{equation}
The first equality is \eqref{eq:weight-square}. For the second,
apply that formula to \(x=z^{*2}\):
\[
 J_rx=(J_rz)^{*2}+(J_{r/2}z)^{\circ2}.
\]
Each \(J_i z\) lies in the basepoint component and is primitive and
\(A\)-annihilated, by Lemma~\ref{lem:weight-product}. Thus
\((J_rz)^{*2}\) is a primitive Pontryagin decomposable, and its
block square is zero by Lemma~\ref{lem:pairing}. Block commutativity
over \(\F_2\) cancels the mixed terms. This gives
\((J_rx)^{\circ2}=(J_{r/2}z)^{\circ4}\).
The indices are integral because \(r\) is a power of two at least two,
and all block products have total weight at most \(s\).

For the last equality in \eqref{eq:power-projection}, observe that
stable counit followed by an adjoint recovers the original spectrum
map.  Thus, on positive-dimensional homology,
\begin{equation}
\label{eq:projected-counit}
 \varepsilon_{i*}J_i=(\operatorname{pr}_i\tau)_*.
\end{equation}
In particular \(B=J_{r/2}z\) has zero stable homology evaluation,
because \(r/2\) is below the leading weight of \(z\). It is a
primitive, \(A\)-annihilated class in \(Q_0X_{r/2}\).
Lemma~\ref{lem:block-nilpotence} gives \(B^{\circ3}=0\), hence
\(B^{\circ4}=0\). This proves the last equality in
\eqref{eq:power-projection}.
The same counit identity gives
\begin{equation}
\label{eq:root-stable-projection}
 \varepsilon_{s*}J_sx=p.
\end{equation}

\emph{Contradictory cone squares.}
Let
\[
 g:S^{2b+1}\longrightarrow T=Q\Sigma X_s
\]
be the adjoint of $j_s\circ f:S^{2b}\to Q_0X_s$, using the
equivalence $QX_s\simeq\Omega Q\Sigma X_s$.
The space $\Sigma X_s$, and hence $T$, is connected.
The loop product under this equivalence is the additive Pontryagin
product.  Equation~\eqref{eq:power-projection} therefore says that
the loop adjoint has Hurewicz class $(J_sx)^{*2}$, and gives
$h(g)=0$ after homology suspension.

Let $\varepsilon:\Sigma^\infty T\to\Sigma^\infty\Sigma X_s$ be
the stable counit and define
\[
 u=\varepsilon^*(\Sigma v)\in H^{b+1}T.
\]
Counit compatibility with homology suspension and
\eqref{eq:root-stable-projection} give
\[
 \langle u,\sigma_*J_sx\rangle
   =\langle v,\varepsilon_{s*}J_sx\rangle
   =\langle v,p\rangle=1.
\]
Stable Steenrod naturality gives
\[
 u^2=\Sq^{b+1}u
 =\varepsilon^*(\Sigma\Sq^{b+1}v)=0.
\]
With $m=b+1$, the hypotheses of Lemma~\ref{lem:cone-square} now hold.
The unique extension $\widetilde u\in H^mC_g$ consequently satisfies
\begin{equation}
\label{eq:nonzero-power-cone-square}
 \Sq^m\widetilde u\ne0.
\end{equation}

On the other hand, the restriction of
$\Sq^{m-1}\widetilde u$ to $T$ is
\[
 \Sq^b u=\varepsilon^*(\Sigma\Sq^b v)
        =\varepsilon^*(\Sigma(v^2)).
\]
This class has an integral torsion lift.  Indeed, for real vector
bundles the reduction of the integral Pontryagin class is
\[
 \rho_2(p_i)=w_{2i}^{\,2}.
\]
For completeness, this follows from
$p_i(\xi)=(-1)^ic_{2i}(\xi\otimes\C)$: the underlying real bundle
of $\xi\otimes\C$ is $\xi\oplus\xi$, and reduction of its Chern
classes gives its even Stiefel--Whitney classes.
Because $v$ is a polynomial in even-index Stiefel--Whitney classes,
its square is the reduction of an integral polynomial
$V\in H^{2b}(B\Sigma_s;\Z)$ in the Pontryagin classes of $\rho_s$.
The degree is positive.  Restriction to the trivial subgroup followed
by transfer shows that positive-degree integral cohomology of the
finite group $\Sigma_s$ is annihilated by $|\Sigma_s|$; thus $V$ is
torsion.  The integral class
\[
 U=\varepsilon^*(\Sigma V)\in H^{2b+1}(T;\Z)
\]
is torsion as well.

Since $H^{2b+1}(S^{2b+1};\Z)=\Z$ is torsion-free, $g^*U=0$.
The integral cohomology sequence of the mapping cone supplies an
extension $\widetilde U\in H^{2b+1}(C_g;\Z)$.
Restriction
\[
 H^{2b+1}(C_g;\F_2)\longrightarrow H^{2b+1}(T;\F_2)
\]
is injective, since the preceding group in the cone sequence is
$\widetilde H^{2b}(S^{2b+1};\F_2)=0$.
Consequently
\[
 \rho_2(\widetilde U)=\Sq^b\widetilde u.
\]
The mod-$2$ Bockstein $\Sq^1$ annihilates an integral reduction.
Finally $b$ is even, so the Adem relation gives
\[
 \Sq^{b+1}\widetilde u
 =\Sq^1\Sq^b\widetilde u
 =\Sq^1\rho_2(\widetilde U)=0,
\]
contradicting \eqref{eq:nonzero-power-cone-square}.
\end{proof}
